% Report for the Style/Mood LoRA project on Stable Diffusion 1.5.
% Build with:  latexmk -pdf report.tex     (or: pdflatex report.tex, twice)
% The figures in figures/ are regenerated from the files in outputs/samples/ and data/processed/.

\documentclass[11pt,a4paper]{article}

\usepackage[T1]{fontenc}
\usepackage[utf8]{inputenc}
\usepackage{lmodern}
\usepackage{microtype}
\usepackage[margin=2.5cm]{geometry}
\usepackage{graphicx}
\usepackage{booktabs}
\usepackage{array}
\usepackage{xcolor}
\usepackage{caption}
\usepackage{enumitem}
\usepackage[hidelinks]{hyperref}

% Colours of the two moods: used only where the respective runs are discussed.
\definecolor{brightHue}{HTML}{A66E15}
\definecolor{scaryHue}{HTML}{456B82}

\newcommand{\tok}[1]{\texttt{\small #1}}
\newcommand{\bright}{\textcolor{brightHue}{\textbf{bright}}}
\newcommand{\scary}{\textcolor{scaryHue}{\textbf{scary}}}

% Width of the image cells in the comparison grids.
\newlength{\thumbw}
\setlength{\thumbw}{3.6cm}
\newcommand{\thumb}[1]{\includegraphics[width=\thumbw]{figures/#1}}

% Left-hand column holding the prompt text, with automatic line breaks.
\newcolumntype{P}{>{\raggedright\arraybackslash}p{2.7cm}}

% File names and flags set in \texttt{} are long and unbreakable: a little extra elasticity
% keeps them from sticking out past the margin.
\emergencystretch=2em

\captionsetup{font=small,labelfont=bf}
\setlength{\parskip}{0.4em}
\setlength{\parindent}{0pt}

\title{\vspace{-1.5cm}Style transfer with LoRA adapters on Stable Diffusion 1.5\\[0.3em]
  \large Project Work in Deep Learning}
\author{Giacomo Antonelli}

\begin{document}
\maketitle

\begin{abstract}
\noindent
Two LoRA adapters are trained on the frozen UNet of Stable Diffusion 1.5 to make it learn an
abstract visual mood from 30 and 35 chosen images, and we measure how well that
style generalises to different prompts. 1,594,368 parameters out of 861M are
updated (0.185\%): each adapter weighs 6.1\,MB against the \textasciitilde3.4\,GB of the full
model. Both moods transfer clearly onto subjects foreign to their respective datasets, with
CLIP deltas that are positive but small ($+0.0380$ and $+0.0195$), consistent with the fact that
image-to-image CLIP similarity is dominated by content rather than style.
\end{abstract}

\section{Goal}

The question is whether a rank-8 LoRA, applied to the attention projections alone of a frozen
UNet, can absorb something as hard to pin down as a mood from a few dozen photographs, and
then apply it to prompts that share nothing with those photographs.

Everything stays frozen except the adapter: the VAE, the text encoder and the 861M parameters of
the UNet. The LoRA matrices are injected into the \texttt{to\_q}, \texttt{to\_k}, \texttt{to\_v}
and \texttt{to\_out.0} projections, that is, where the text conditioning enters. The style is
bound to a trigger phrase that ends every training caption, and the same phrase reactivates it at
inference time.

\section{Method}

\subsection{Data}

I've created one subfolder for each mood. Images are brought to $512\times512$ by resizing on the short side and
centre cropping, and each is paired with a caption ending in its own mood's trigger. The trigger
is recorded in \texttt{dataset.json} and propagated all the way to the checkpoint, so it does not
have to be retyped at inference time: a suffix different from the training one leaves the adapter
silently inactive, and this was the most frequent mistake in the early experiments.

\begin{table}[h]
\centering
\begin{tabular}{llrl}
\toprule
run & trigger & images & character of the dataset \\
\midrule
\bright & \tok{<cjh-bright>} & 30 & saturated, luminous, illustrative \\
\scary  & \tok{<cjh-scary>}  & 35 & desaturated, cold, foggy \\
\bottomrule
\end{tabular}
\caption{The two style datasets. The trigger is the only part shared by every caption of a mood.}
\end{table}

\begin{figure}[h]
\centering
\includegraphics[width=0.86\textwidth]{figures/dataset_bright.jpg}\\[0.35em]
\includegraphics[width=0.86\textwidth]{figures/dataset_scary.jpg}
\caption{Three training images per mood: \bright{} on top, \scary{} at the bottom.}
\end{figure}

\subsection{Training}

Both runs share the same hyperparameters, so the two adapters differ only in what they have seen:
rank 8, $\alpha=16$, learning rate $10^{-4}$, 1200 steps, batch 1 with gradient accumulation 2,
resolution $512\times512$, seed 42.

Each step samples a random timestep, adds the corresponding noise to the latent produced by the
VAE and asks the UNet to predict it, the loss is the MSE between predicted and true noise, and the
gradients reach the LoRA matrices only. Gradient checkpointing and fp16 are what make the training
fit within the 6\,GB of the GPU used (GTX 1060), on this Pascal architecture, fp16 saves memory more than it
speeds things up.

\begin{table}[h]
\centering
\begin{tabular}{lr}
\toprule
trainable parameters & 1,594,368 \\
total UNet parameters & 861,115,332 \\
fraction trained & 0.185\,\% \\
size of the adapter & 6.1\,MB \\
size of the full model & \textasciitilde3.4\,GB \\
GPU & GTX 1060 6\,GB \\
duration of one run & 1h\,27m \\
\bottomrule
\end{tabular}
\caption{Cost of the fine-tuning.}
\end{table}

\subsection{Evaluation}

The 11 evaluation prompts describe quite different subjects: four neutral with respect
to both moods, three ambiguous, and two per mood chosen to pull the result one way. The seed
depends only on the prompt index and not on the variant, so base and adapter start from the same
noise and the differences are attributable to the adapter alone.

The metric is the mean CLIP similarity (\texttt{clip-vit-base-patch32}) between the generated
images and the mood's reference dataset. The absolute value is not interpretable on its own, so
what gets reported is the delta against the same images generated by the base model.

\section{Results}

\begin{table}[h]
\centering
\begin{tabular}{lrrrrr}
\toprule
run & base & lora & delta & std & prompts improved \\
\midrule
\bright & 0.5877 & 0.6257 & $+0.0380$ & 0.0673 & 7 / 11 \\
\scary  & 0.5571 & 0.5766 & $+0.0195$ & 0.0643 & 8 / 11 \\
\bottomrule
\end{tabular}
\caption{Mean CLIP similarity against the reference dataset of each mood.}
\end{table}

Both deltas are positive and both are small, which is the expected outcome and not a disappointing
result: image-to-image CLIP similarity is dominated by the subject, and the subject is held fixed
across the columns by construction. The metric should be read as weak evidence; the comparison
grid remains the primary evaluation.

On visual inspection the transfer is already clear at scale 1.0, so no sweep over the adapter
strength was needed. \bright{} turns the photographs into illustration: the lighthouse acquires
painted rock and a poster-like sky, the fisherman becomes a painterly portrait, the waterfall
shifts towards turquoise and violet. \scary{}, by contrast, works by subtraction: it removes
saturation, cools the palette and fills the middle ground with haze, to the point that the forest
turns yellow-grey and the wolf in the snow loses almost all of its colour. It is worth noting that
the scary result show just a change in the mood, rather than a complete transformation (like adding blood, or horror subjects).

Focusing on the rows of the neutral subjects, we can see that even if those are have nothing to do with
 either dataset, they are still restyled coherently, showing a generalizationism of these modules.

The main problem I can see in this project is that the trigger carries meaning: the CLIP tokenizer breaks the trigger into
ordinary pieces \tok{['in','<','cj','h','-','scary','>','style']} and \emph{scary} and
\emph{bright} are English words the base model already understands and might influence the model's behavior.

\clearpage
\appendix
\section{Comparison grids}

Every row uses the same seed across the three columns: the differences are attributable to the
adapter alone. The adapters are applied at scale 1.0.

\begin{figure}[h]
\centering
\begin{tabular}{P ccc}
& \small\textbf{BASE} & \small\textcolor{brightHue}{\textbf{BRIGHT}} & \small\textcolor{scaryHue}{\textbf{SCARY}} \\
\small a lighthouse on a rocky cliff & \thumb{base_00.jpg} & \thumb{bright_00.jpg} & \thumb{scary_00.jpg} \\
\small a portrait of an old fisherman & \thumb{base_01.jpg} & \thumb{bright_01.jpg} & \thumb{scary_01.jpg} \\
\small a bowl of fruit on a wooden table & \thumb{base_02.jpg} & \thumb{bright_02.jpg} & \thumb{scary_02.jpg} \\
\small a city skyline seen from a rooftop & \thumb{base_03.jpg} & \thumb{bright_03.jpg} & \thumb{scary_03.jpg} \\
\end{tabular}
\caption{Subjects neutral with respect to both datasets: they measure how far the style applies
outside its own domain.}
\end{figure}

\clearpage

\begin{figure}[h]
\centering
\begin{tabular}{P ccc}
& \small\textbf{BASE} & \small\textcolor{brightHue}{\textbf{BRIGHT}} & \small\textcolor{scaryHue}{\textbf{SCARY}} \\
\small a forest & \thumb{base_04.jpg} & \thumb{bright_04.jpg} & \thumb{scary_04.jpg} \\
\small a wolf standing in a snowy forest & \thumb{base_05.jpg} & \thumb{bright_05.jpg} & \thumb{scary_05.jpg} \\
\small an empty subway station at night & \thumb{base_06.jpg} & \thumb{bright_06.jpg} & \thumb{scary_06.jpg} \\
\small an abandoned hospital corridor & \thumb{base_07.jpg} & \thumb{bright_07.jpg} & \thumb{scary_07.jpg} \\
\end{tabular}
\caption{Ambiguous subjects: it is the adapter's mood that decides which way they lean.}
\end{figure}

\clearpage

\begin{figure}[h]
\centering
\begin{tabular}{P ccc}
& \small\textbf{BASE} & \small\textcolor{brightHue}{\textbf{BRIGHT}} & \small\textcolor{scaryHue}{\textbf{SCARY}} \\
\small a scarecrow in a foggy field at dusk & \thumb{base_08.jpg} & \thumb{bright_08.jpg} & \thumb{scary_08.jpg} \\
\small a meadow of wildflowers under a clear sky & \thumb{base_09.jpg} & \thumb{bright_09.jpg} & \thumb{scary_09.jpg} \\
\small a waterfall in a sunlit tropical forest & \thumb{base_10.jpg} & \thumb{bright_10.jpg} & \thumb{scary_10.jpg} \\
\end{tabular}
\caption{Subjects that push towards one mood or the other: the first towards \scary{}, the other
two towards \bright{}.}
\end{figure}

\end{document}
